% GENERATED FILE. Edit canonical lessons, then run npm run generate:books.
% canonical-source-hash: fnv1a64:940d6cd9f774c8a5
% canonical-lessons: SA-C303-durbhiksam, SA-C303-mahamari, SA-C303-ksinah, SA-C303-vyapnoti, SA-C303-mlayati

\chapter{Famine and Sickness}
\label{ch:sa-famine-and-sickness}
\hlchaptermodality{303}

\begin{chapteropening}
\textbf{By the end of this chapter:} \emph{I can tell how hunger and disease spread.}

Complete the last lesson of chapter 303: I can tell how hunger and disease spread.
\end{chapteropening}

\section[\texorpdfstring{\sk{दुर्भिक्षम्} (durbhikṣam)}{दुर्भिक्षम् (durbhikṣam)}]{\sk{दुर्भिक्षम्} (durbhikṣam) — a famine}

\label{lesson:SA-C303-durbhiksam}



\begin{quote}
Before the new one: say the Sanskrit for goes under, then the Sanskrit for by degrees.
\end{quote}

\subsection*{You'll want to know: \sk{दुर्भिक्षम्}}
\textbf{\sk{दुर्भिक्षम्}} — \emph{durbhikṣam} — \textquotedblleft{}a famine\textquotedblright{}.

\begin{quote}
\sk{तस्मिन्} \sk{संवत्सरे} \sk{दुर्भिक्षम्} \sk{अभवत्।}
\end{quote}

\begin{quote}
\emph{tasmin saṃvatsare durbhikṣam abhavat.}
\end{quote}

\begin{quote}
That year there was a famine.
\end{quote}

A neuter noun.

\begin{cousinweb}
\textbf{\sk{दुर्भिक्षम्}} (\emph{durbhikṣam}) is built from \emph{dur}, \textquotedblleft{}bad, hard\textquotedblright{}, and \emph{bhikṣā}, \textquotedblleft{}alms, food begged\textquotedblright{}; MW gives \emph{durbhikṣa} as scarcity of provisions, famine (MW).
\end{cousinweb}

\begin{grammarlens}[title={What you've built}]
One more word for telling what happened.
\end{grammarlens}

\subsection*{Your turn}
\begin{itemize}
\raggedright
  \item \emph{Say it:} \emph{durbhikṣam}
  \item \emph{Say it:} \emph{durbhikṣam}, once more
  \item \emph{Say it:} \emph{tasmin saṃvatsare durbhikṣam abhavat.} — That year there was a famine.
  \item \emph{Recall:} say the Sanskrit for goes under, then the Sanskrit for by degrees, then say \emph{durbhikṣam} again
\end{itemize}

\begin{tcolorbox}[breakable,colback=teal!4,colframe=teal!35!black,title={Before you move on}]
What does \sk{दुर्भिक्षम्} mean? (\textquotedblleft{}A famine\textquotedblright{}.) Say it once more.
\end{tcolorbox}
\section[\texorpdfstring{\sk{महामारी} (mahāmārī)}{महामारी (mahāmārī)}]{\sk{महामारी} (mahāmārī) — an epidemic}

\label{lesson:SA-C303-mahamari}



\begin{quote}
Before the new one: say the Sanskrit for by degrees, then the Sanskrit for a famine.
\end{quote}

\subsection*{You'll want to know: \sk{महामारी}}
\textbf{\sk{महामारी}} — \emph{mahāmārī} — \textquotedblleft{}an epidemic\textquotedblright{}.

\begin{quote}
\sk{नगरे} \sk{महामारी} \sk{प्रसरति।}
\end{quote}

\begin{quote}
\emph{nagare mahāmārī prasarati.}
\end{quote}

\begin{quote}
An epidemic is spreading in the city.
\end{quote}

A feminine noun.

\begin{cousinweb}
\textbf{\sk{महामारी}} (\emph{mahāmārī}) is built from \emph{mahā}, \textquotedblleft{}great\textquotedblright{}, and \emph{mārī}, \textquotedblleft{}pestilence\textquotedblright{}; MW gives \emph{mahāmārī} as a pestilence that kills many, cholera (MW).
\end{cousinweb}

\begin{grammarlens}[title={What you've built}]
One more word for telling what happened.
\end{grammarlens}

\subsection*{Your turn}
\begin{itemize}
\raggedright
  \item \emph{Say it:} \emph{mahāmārī}
  \item \emph{Say it:} \emph{mahāmārī}, once more
  \item \emph{Say it:} \emph{nagare mahāmārī prasarati.} — An epidemic is spreading in the city.
  \item \emph{Recall:} say the Sanskrit for by degrees, then the Sanskrit for a famine, then say \emph{mahāmārī} again
\end{itemize}

\begin{tcolorbox}[breakable,colback=teal!4,colframe=teal!35!black,title={Before you move on}]
What does \sk{महामारी} mean? (\textquotedblleft{}An epidemic\textquotedblright{}.) Say it once more.
\end{tcolorbox}
\section[\texorpdfstring{\sk{क्षीणः} (kṣīṇaḥ)}{क्षीणः (kṣīṇaḥ)}]{\sk{क्षीणः} (kṣīṇaḥ) — wasted away; weakened}

\label{lesson:SA-C303-ksinah}



\begin{quote}
Before the new one: say the Sanskrit for a famine, then the Sanskrit for an epidemic.
\end{quote}

\subsection*{You'll want to know: \sk{क्षीणः}}
\textbf{\sk{क्षीणः}} — \emph{kṣīṇaḥ} — \textquotedblleft{}wasted away; weakened\textquotedblright{}.

\begin{quote}
\sk{क्षुधया} \sk{बालाः} \sk{क्षीणाः} \sk{सन्ति।}
\end{quote}

\begin{quote}
\emph{kṣudhayā bālāḥ kṣīṇāḥ santi.}
\end{quote}

\begin{quote}
The children are wasted with hunger.
\end{quote}

\begin{cousinweb}
\textbf{\sk{क्षीणः}} (\emph{kṣīṇaḥ}) comes from the root √kṣi, \textquotedblleft{}to waste, destroy\textquotedblright{}; MW gives \emph{kṣīṇa} as diminished, wasted, worn away, emaciated (MW).
\end{cousinweb}

\begin{grammarlens}[title={What you've built}]
One more word for telling what happened.
\end{grammarlens}

\subsection*{Your turn}
\begin{itemize}
\raggedright
  \item \emph{Say it:} \emph{kṣīṇaḥ}
  \item \emph{Say it:} \emph{kṣīṇaḥ}, once more
  \item \emph{Say it:} \emph{kṣudhayā bālāḥ kṣīṇāḥ santi.} — The children are wasted with hunger.
  \item \emph{Recall:} say the Sanskrit for a famine, then the Sanskrit for an epidemic, then say \emph{kṣīṇaḥ} again
\end{itemize}

\begin{tcolorbox}[breakable,colback=teal!4,colframe=teal!35!black,title={Before you move on}]
What does \sk{क्षीणः} mean? (\textquotedblleft{}Wasted away; weakened\textquotedblright{}.) Say it once more.
\end{tcolorbox}
\section[\texorpdfstring{\sk{व्याप्नोति} (vyāpnoti)}{व्याप्नोति (vyāpnoti)}]{\sk{व्याप्नोति} (vyāpnoti) — spreads through; fills}

\label{lesson:SA-C303-vyapnoti}



\begin{quote}
Before the new one: say the Sanskrit for an epidemic, then the Sanskrit for wasted away.
\end{quote}

\subsection*{You'll want to know: \sk{व्याप्नोति}}
\textbf{\sk{व्याप्नोति}} — \emph{vyāpnoti} — \textquotedblleft{}spreads through; fills\textquotedblright{}.

\begin{quote}
\sk{रोगः} \sk{सर्वं} \sk{नगरं} \sk{व्याप्नोति।}
\end{quote}

\begin{quote}
\emph{rogaḥ sarvaṃ nagaraṃ vyāpnoti.}
\end{quote}

\begin{quote}
The disease spreads through the whole city.
\end{quote}

Past (it …): \sk{रोगः} \sk{नगरं} \sk{व्याप्तवान्}, \emph{rogaḥ nagaraṃ vyāptavān}.

\begin{cousinweb}
\textbf{\sk{व्याप्नोति}} (\emph{vyāpnoti}) is built from \emph{vi} and the root √āp, \textquotedblleft{}to reach\textquotedblright{}, which also gives \emph{prāpnoti}; MW gives \emph{vy-āp} as to reach through, pervade, fill (MW).
\end{cousinweb}

\begin{grammarlens}[title={What you've built}]
One more word for telling what happened.
\end{grammarlens}

\subsection*{Your turn}
\begin{itemize}
\raggedright
  \item \emph{Say it:} \emph{vyāpnoti}
  \item \emph{Say it:} \emph{vyāpnoti}, once more
  \item \emph{Say it:} \emph{rogaḥ sarvaṃ nagaraṃ vyāpnoti.} — The disease spreads through the whole city.
  \item \emph{Recall:} say the Sanskrit for an epidemic, then the Sanskrit for wasted away, then say \emph{vyāpnoti} again
\end{itemize}

\begin{tcolorbox}[breakable,colback=teal!4,colframe=teal!35!black,title={Before you move on}]
What does \sk{व्याप्नोति} mean? (\textquotedblleft{}Spreads through; fills\textquotedblright{}.) Say it once more.
\end{tcolorbox}
\section[\texorpdfstring{\sk{म्लायति} (mlāyati)}{म्लायति (mlāyati)}]{\sk{म्लायति} (mlāyati) — withers; fades}

\label{lesson:SA-C303-mlayati}



\begin{quote}
Before the new one: say the Sanskrit for wasted away, then the Sanskrit for spreads through.
\end{quote}

\subsection*{You'll want to know: \sk{म्लायति}}
\textbf{\sk{म्लायति}} — \emph{mlāyati} — \textquotedblleft{}withers; fades\textquotedblright{}.

\begin{quote}
\sk{वर्षां} \sk{विना} \sk{पुष्पाणि} \sk{म्लायन्ति।}
\end{quote}

\begin{quote}
\emph{varṣāṃ vinā puṣpāṇi mlāyanti.}
\end{quote}

\begin{quote}
Without rain the flowers wither.
\end{quote}

Past (it …): \sk{पुष्पं} \sk{म्लानवत्}, \emph{puṣpaṃ mlānavat}.

\begin{cousinweb}
\textbf{\sk{म्लायति}} (\emph{mlāyati}) comes from the root √mlai, \textquotedblleft{}to fade, wither, decay\textquotedblright{}, which MW also gives as to be languid or exhausted (MW).
\end{cousinweb}

\begin{grammarlens}[title={What you've built}]
One more word for telling what happened.
\end{grammarlens}

\subsection*{Your turn}
\begin{itemize}
\raggedright
  \item \emph{Say it:} \emph{mlāyati}
  \item \emph{Say it:} \emph{mlāyati}, once more
  \item \emph{Say it:} \emph{varṣāṃ vinā puṣpāṇi mlāyanti.} — Without rain the flowers wither.
  \item \emph{Recall:} say the Sanskrit for wasted away, then the Sanskrit for spreads through, then say \emph{mlāyati} again
\end{itemize}

\begin{tcolorbox}[breakable,colback=teal!4,colframe=teal!35!black,title={Before you move on}]
What does \sk{म्लायति} mean? (\textquotedblleft{}Withers; fades\textquotedblright{}.) Say it once more.
\end{tcolorbox}
